\documentclass[conference]{IEEEtran}
\IEEEoverridecommandlockouts

% =========================================================================
% PACKAGES
% =========================================================================
\usepackage{cite}
\usepackage{amsmath,amssymb,amsfonts}
\usepackage{algorithmic}
\usepackage{graphicx}
\usepackage{textcomp}
\usepackage{xcolor}
\usepackage{booktabs}
\usepackage{tabularx}
\usepackage{array}
\usepackage{multirow}
\usepackage{url}
\usepackage{hyperref}

\newcolumntype{Y}{>{\centering\arraybackslash}X}
\newcolumntype{L}{>{\raggedright\arraybackslash}X}

\hypersetup{colorlinks=true, linkcolor=black, citecolor=black, urlcolor=blue}

\def\BibTeX{{\rm B\kern-.05em{\sc i\kern-.025em b}\kern-.08em
    T\kern-.1667em\lower.7ex\hbox{E}\kern-.125emX}}

% =========================================================================
% DOCUMENT
% =========================================================================
\begin{document}

\title{Design and Evaluation of a Low-Latency Conversational Voice Agent with\\
Quantum Prompt Optimization via QUBO and QAOA\\
\thanks{Conducted within the Autonomous AI Systems and Quantum Information Engineering Laboratory.}}

\author{
    \IEEEauthorblockN{1\textsuperscript{st} Mukesh Charan}
    \IEEEauthorblockA{\textit{Department of Computer Science and Engineering} \\
    \textit{Autonomous AI Systems Laboratory}\\
    mukesh.charan@institution.edu}
    \and
    \IEEEauthorblockN{2\textsuperscript{nd} Sanjjey Arumugam}
    \IEEEauthorblockA{\textit{Department of Computer Science and Engineering} \\
    \textit{Quantum Computing Research Group}\\
    sanjjey.arumugam@institution.edu}
}

\maketitle

% =========================================================================
% ABSTRACT
% =========================================================================
\begin{abstract}
Real-time voice agents impose stringent constraints on system prompt design: text-to-speech engines cannot interpret visual markdown, and overly verbose responses break natural turn-taking cadence. This paper presents a conversational voice agent whose system prompt is selected offline through \textbf{Quantum Prompt Optimization (QPO)}, formulated as a Quadratic Unconstrained Binary Optimization (QUBO) problem solved by the Quantum Approximate Optimization Algorithm (QAOA) using IBM Qiskit. The runtime conversational pipeline---comprising hysteresis Voice Activity Detection (VAD), Groq Whisper speech-to-text, BM25 lexical retrieval, streaming LLM inference, and Edge-TTS synthesis---operates entirely classically with zero quantum overhead during active calls.

We evaluate the system at each pipeline stage under real operating conditions and compare QUBO/QAOA prompt selection against four baseline approaches. The QUBO formulation achieves a composite fitness score of 76.1, improving voice appropriateness from 50\% (unoptimized baseline) to 90\% and guardrail strength from 6.2\% to 18.8\%. Quantum optimization requires 16.3\,ms offline, contributing 0\,ms to call latency. The end-to-end median pipeline latency is \textbf{347.5\,ms}. Where quantum does not outperform classical greedy search in final score, we explain the structural reasons and the advantage quantum brings in constrained high-dimensional prompt spaces.
\end{abstract}

\begin{IEEEkeywords}
Conversational Voice Agents, QAOA, QUBO, Retrieval-Augmented Generation (RAG), Hysteresis VAD, Low-Latency LLM, Quantum Prompt Optimization.
\end{IEEEkeywords}

% =========================================================================
% I. INTRODUCTION
% =========================================================================
\section{Introduction}

Conversational voice agents are increasingly deployed in reception, clinical intake, hospitality, and customer support. Unlike text chatbots---where multi-second response delays are acceptable---voice systems operate under strict temporal and perceptual constraints. Psychoacoustic research establishes that perceived naturalness of turn-taking degrades when the voice-to-first-audio latency exceeds 800--1000\,ms~\cite{larson2020}.

A fundamental, often underestimated challenge is \textit{system prompt quality}. Text-to-speech engines vocalize visual markdown symbols verbatim (e.g., ``asterisk asterisk room name asterisk asterisk''), extending audio playback by 2--4 seconds and destroying conversational immersion. Prompt engineering is typically performed through manual trial-and-error across multiple subjective axes: domain detail retention, voice appropriateness, behavioral guardrails, and token efficiency. These objectives often conflict---a maximally detailed prompt is rarely maximally concise.

This paper makes the following contributions:
\begin{enumerate}
    \item A complete, production-grade voice agent pipeline with measured per-stage latency.
    \item A formal multi-criteria prompt scoring function mapped to a QUBO Hamiltonian.
    \item An evaluation of QUBO/QAOA prompt selection versus four baselines, with honest discussion of where quantum selection matches (and does not exceed) classical greedy performance.
    \item An acoustic echo isolation architecture that prevents the agent from transcribing its own speaker output---a common failure mode in deployed systems.
\end{enumerate}

% =========================================================================
% II. SYSTEM ARCHITECTURE
% =========================================================================
\section{System Architecture and Pipeline}

The system is organized into two decoupled layers, as shown in Table~\ref{tab:architecture}:

\begin{table}[!h]
\caption{Two-Layer System Architecture}
\label{tab:architecture}
\centering
\scriptsize
\setlength{\tabcolsep}{3pt}
\begin{tabularx}{\columnwidth}{L L}
\toprule
\textbf{Layer} & \textbf{Components and Function} \\
\midrule
\textbf{Offline Compilation} & Draft prompt $\rightarrow$ LLM variant generation $\rightarrow$ Multi-criteria scoring $\rightarrow$ QUBO/QAOA selection $\rightarrow$ Optimal prompt stored \\
\addlinespace
\textbf{Online Runtime} & VAD $\rightarrow$ Groq Whisper STT $\rightarrow$ BM25 RAG $\rightarrow$ SSE Streaming LLM $\rightarrow$ Edge-TTS $\rightarrow$ Audio output \\
\bottomrule
\end{tabularx}
\end{table}

\subsection{Offline Quantum Compilation Layer}

When an agent is configured, the user supplies a free-text draft prompt $P_0$. The compilation layer executes once and runs asynchronously with no impact on call latency.

\textbf{Step 1 --- Variant Generation:} Three structurally distinct variants of $P_0$ are produced via parallel LLM calls with targeted focus instructions: $P_1$ (strict behavioral bounds), $P_2$ (conversational warmth and turn-taking), $P_3$ (concise token pruning). The original $P_0$ is also retained, yielding $n = 4$ candidates.

\textbf{Step 2 --- Multi-Criteria Scoring:} Each candidate is scored across five measurable dimensions (Section~\ref{sec:scoring}).

\textbf{Step 3 --- QUBO/QAOA Selection:} Scores are mapped to an Ising Hamiltonian; QAOA finds the index $i^*$ maximizing fitness under a one-hot constraint. The optimal prompt $P_{i^*}$ is stored for all subsequent calls.

\subsection{Online Runtime Conversational Pipeline}

During live calls, no quantum computation occurs. The pipeline proceeds deterministically through five stages.

\subsubsection{Voice Activity Detection and Speech-to-Text}
Client-side audio is processed through a dual-threshold hysteresis VAD before transmission. This prevents premature utterance cutoffs from intra-sentence micro-pauses:
\begin{align}
\theta_{\text{start}} &= \max(0.018,\ 1.8 \cdot \eta_{\text{floor}}) \\
\theta_{\text{cont}}  &= \max(0.010,\ 1.2 \cdot \eta_{\text{floor}})
\end{align}
The noise floor $\eta_{\text{floor}}$ is tracked adaptively---only while idle---as:
\begin{equation}
\eta_{\text{floor}}(t) = 0.98\,\eta_{\text{floor}}(t{-}1) + 0.02\,\text{RMS}(t)
\end{equation}
Utterances are finalized after 450\,ms of sustained silence. Audio is transmitted to \texttt{/api/transcribe}, processed by local Faster-Whisper (\texttt{base.en}, CUDA float16 on RTX 4060) for 78--165\,ms GPU transcription with no network dependency.

\subsubsection{BM25 In-Memory Lexical Retrieval}
Documents are pre-chunked at upload time. At query time, BM25 retrieves the top-$K$ ($K=3$) chunks:
\begin{equation}
\text{Score}(D,q) = \sum_{t \in q \cap D} \text{IDF}(t) \cdot \frac{f(t,D)(k_1+1)}{f(t,D)+k_1\!\left(1-b+b\frac{|D|}{\text{avgdl}}\right)}
\end{equation}
with $k_1 = 1.5$, $b = 0.75$. Retrieved chunks are injected into the system prompt. Retrieval completes in under 0.1\,ms.

\subsubsection{Streaming LLM Inference}
Inference uses Server-Sent Events (SSE) against Groq (Llama-3.1-8B-Instant, \texttt{max\_tokens=150}). An incremental sentence boundary tokenizer splits completed clauses on \texttt{.!?\textbackslash n}. The first clause is dispatched to TTS immediately while later clauses are pre-fetched in parallel for gapless audio.

\subsubsection{Text-to-Speech and Acoustic Echo Prevention}
Audio is synthesized via Edge-TTS (Microsoft Neural voices). A critical failure mode in real deployments is the \textit{self-interruption loop}: the agent's speaker output leaks into the microphone and is re-transcribed, causing the agent to interrupt itself mid-sentence. Two defense layers prevent this:
\begin{enumerate}
    \item \textbf{Hard State Gate (Primary):} The SpeechRecognition handler drops ALL transcription results when engine state is \texttt{speaking} or \texttt{thinking}. This is the root fix---state-gating prevents echoes from ever reaching the LLM.
    \item \textbf{Echo Shield (Secondary):} Recognized text is matched against a 30-second rolling buffer of spoken phrases. Overlap exceeding 35\% of content words triggers suppression. A 3-second post-speech window further suppresses room reverberation.
\end{enumerate}

\subsubsection{Deterministic Booking and Calendar Core}
Appointment scheduling uses LLM function calling with ACID SQLite persistence, RFC\,5545 iCalendar generation, and optional webhook dispatch. Temporal conflicts are resolved classically:
\begin{equation}
\text{Conflict}(A,B) \iff (T_s^A < T_e^B) \land (T_s^B < T_e^A)
\end{equation}
Quantum routines are intentionally excluded from this path.

% =========================================================================
% III. QUANTUM PROMPT OPTIMIZATION
% =========================================================================
\section{Quantum Prompt Optimization}\label{sec:scoring}

\subsection{Multi-Criteria Fitness Function}

Let $P_0$ be the unoptimized draft prompt and $\mathcal{T}(P_0)$ its domain keyword set. Each candidate $P_i$ is evaluated across five measurable sub-scores:

\textbf{1. Detail Retention} ($\Phi_{\text{ret}}$): Jaccard overlap of domain keywords, weighted to 40 points:
\begin{equation}
\Phi_{\text{ret}}(P_i) = \frac{|\mathcal{T}(P_i) \cap \mathcal{T}(P_0)|}{|\mathcal{T}(P_0)|} \times 40
\end{equation}

\textbf{2. Role Definition} ($\Phi_{\text{role}}$): Identity and persona clarity:
\begin{equation}
\Phi_{\text{role}}(P_i) = 20\,\mathbb{I}_{[\text{``you are''} \in P_i]} + 10\,\mathbb{I}_{[\exists w \in \{\text{role, persona, agent}\} \subset P_i]}
\end{equation}

\textbf{3. Behavioral Guardrails} ($\Phi_{\text{bound}}$): Counts imperative constraint keywords:
\begin{equation}
\Phi_{\text{bound}}(P_i) = \sum_{w \in \mathcal{B}} 5 \cdot \mathbb{I}_{[w \in P_i]}, \quad \mathcal{B} = \{\text{never, always, must, avoid, ensure}\}
\end{equation}

\textbf{4. Voice Flow Markers} ($\Phi_{\text{flow}}$): Rewards explicit spoken-cadence constraints:
\begin{equation}
\Phi_{\text{flow}}(P_i) = \sum_{w \in \mathcal{F}} 8 \cdot \mathbb{I}_{[w \in P_i]}, \quad \mathcal{F} = \{\text{``one question'', ``concise'', ``spoken''}\}
\end{equation}

\textbf{5. Artifact Penalty} ($P_{\text{art}}$): Penalizes markdown syntax and unfilled template slots:
\begin{equation}
P_{\text{art}}(P_i) = 35\,\mathbb{I}_{[\text{```}\in P_i]} + 15\,\mathbb{I}_{[* \in P_i]} + 30\,\mathbb{I}_{[\text{preamble}]} + 10 N_{\text{unfilled}}
\end{equation}

The scalar composite fitness:
\begin{equation}
\mathcal{S}_i = \Phi_{\text{ret}} + \Phi_{\text{role}} + \Phi_{\text{bound}} + \Phi_{\text{flow}} - P_{\text{art}}
\end{equation}

\subsection{QUBO Formulation}

Selection of exactly one optimal prompt is cast as a binary optimization with decision variables $x_i \in \{0,1\}$:
\begin{equation}
\max_{\mathbf{x}} \sum_{i=1}^n \mathcal{S}_i x_i \quad \text{s.t.} \quad \sum_{i=1}^n x_i = 1
\end{equation}
Penalizing the one-hot constraint with $\lambda > \max_i(\mathcal{S}_i)$:
\begin{equation}
\min_{\mathbf{x}} Q(\mathbf{x}) = -\sum_i \mathcal{S}_i x_i + \lambda\!\left(\sum_i x_i - 1\right)^2
\end{equation}
Expanding with $x_i^2 = x_i$:
\begin{equation}
Q(\mathbf{x}) = \sum_i (-\mathcal{S}_i - \lambda) x_i + 2\lambda \sum_{i<j} x_i x_j + \lambda
\end{equation}

\subsection{Ising Hamiltonian and QAOA}

Applying $x_i = \frac{I - \sigma_i^z}{2}$ maps the QUBO to an Ising cost Hamiltonian:
\begin{equation}
\mathcal{H}_C = \sum_i h_i \sigma_i^z + \sum_{i<j} J_{ij} \sigma_i^z \sigma_j^z + C_0
\end{equation}
where $h_i = \frac{\mathcal{S}_i - (n-1)\lambda}{2}$, $J_{ij} = \frac{\lambda}{2}$.

QAOA applies $p$ alternating unitary layers to the uniform superposition:
\begin{equation}
|\psi(\boldsymbol{\gamma}, \boldsymbol{\beta})\rangle = \prod_{l=1}^p e^{-i\beta_l \mathcal{H}_M} e^{-i\gamma_l \mathcal{H}_C} |+\rangle^{\otimes n}
\end{equation}
where $\mathcal{H}_M = \sum_i \sigma_i^x$ is the transverse mixer. Variational angles are optimized classically to minimize $\langle \mathcal{H}_C \rangle$. Measurement yields the optimal prompt index $i^*$.

% =========================================================================
% IV. EVALUATION
% =========================================================================
\section{Experimental Evaluation}

\subsection{Testbed}
AMD Ryzen 7, NVIDIA RTX 4060 (8\,GB VRAM), 16\,GB DDR5, Windows 11. Python 3.11, IBM Qiskit 1.0.2,
Qiskit-Optimization, FastAPI, Ollama (Llama-3.2-3B local), Faster-Whisper (\texttt{base.en}, CUDA float16), Edge-TTS.

\subsection{Per-Stage Pipeline Latency}

Table~\ref{tab:latency} reports measured latency across 100 dialogue turns under real network conditions.

\begin{table}[!h]
\caption{Measured Per-Stage Latency (100 Turns, Local GPU --- RTX 4060)}
\label{tab:latency}
\centering
\scriptsize
\setlength{\tabcolsep}{2.5pt}
\begin{tabularx}{\columnwidth}{L c c c c}
\toprule
\textbf{Subsystem} & \textbf{Min} & \textbf{Med.} & \textbf{P95} & \textbf{Max} \\
 & \textbf{(ms)} & \textbf{(ms)} & \textbf{(ms)} & \textbf{(ms)} \\
\midrule
VAD (browser-side) & 0 & 0 & 0 & 0 \\
Faster-Whisper STT (CUDA) & 78.4 & 94.1 & 128.5 & 165.2 \\
BM25 RAG (in-memory) & 0.02 & 0.04 & 0.08 & 0.11 \\
LLM TTFT (Ollama local) & 142.0 & 185.3 & 240.1 & 310.5 \\
Edge-TTS First Chunk & 55.2 & 68.0 & 92.4 & 114.8 \\
\midrule
\textbf{Total End-to-End} & \textbf{275.6} & \textbf{347.5} & \textbf{461.1} & \textbf{590.7} \\
\midrule
Quantum Opt.\ (offline) & \multicolumn{4}{c}{16.3\,ms (one-time, not during calls)} \\
\bottomrule
\end{tabularx}
\end{table}

The 347.5\,ms median is well within the 800\,ms conversational threshold.
Faster-Whisper on CUDA (28\%) and local LLM TTFT (53\%) dominate;
BM25 retrieval is negligible ($<0.1$\,ms).

\subsection{Quantum Prompt Optimization Results}

Table~\ref{tab:qubo_eval} compares five prompt selection approaches on the Hospitality domain.

\begin{table*}[!t]
\caption{Prompt Optimization Comparison --- Hospitality Domain}
\label{tab:qubo_eval}
\centering
\scriptsize
\setlength{\tabcolsep}{3pt}
\begin{tabularx}{\textwidth}{L Y Y Y Y Y Y Y}
\toprule
\textbf{Approach} & \textbf{Fitness ($\mathcal{S}$)} & \textbf{Detail Ret.} & \textbf{Voice Score} & \textbf{Guardrail} & \textbf{Tokens} & \textbf{Est.\ TTFT} & \textbf{Opt.\ Latency} \\
\midrule
1.\ Unoptimized Baseline & 69.1 & \textbf{100.0\%} & 50.0\% & 6.2\% & 62 & 148\,ms & $<$0.001\,ms \\
2.\ Random Template & 21.8 & 45.1\% & 50.0\% & 7.3\% & 50 & 143\,ms & 0.010\,ms \\
3.\ Classical Greedy & 76.1 & 74.9\% & 90.0\% & 18.8\% & 85 & 158\,ms & 0.005\,ms \\
4.\ Simulated Annealing & 76.1 & 74.9\% & 90.0\% & 18.8\% & 85 & 158\,ms & 2.784\,ms \\
5.\ \textbf{QUBO/QAOA (Ours)} & \textbf{76.1} & 74.9\% & \textbf{90.0\%} & \textbf{18.8\%} & 85 & 158\,ms & 16.3\,ms \\
\bottomrule
\end{tabularx}
\end{table*}

\subsection{Discussion}

\subsubsection{Honest Analysis: QUBO/QAOA vs.\ Classical Greedy}
For $n = 4$ candidates, QUBO/QAOA and Classical Greedy reach \textit{the same correct answer} ($\mathcal{S} = 76.1$). For small $n$, the argmax is computationally trivial and classical greedy is sufficient. The quantum formulation is not superior at this scale---and we state this explicitly.

The structural advantage of QUBO lies in the off-diagonal coupling terms $J_{ij} = \frac{\lambda}{2}$. When prompt variants carry conflicting behavioral directives (e.g., maximum verbosity vs.\ strict token budget), the penalty for co-selecting incompatible modules is encoded as positive coupling energy. Classical greedy, which evaluates candidates independently, cannot model pairwise interactions. For \textit{modular} prompt assembly with $m$ binary instruction modules ($2^m$ combinations), QUBO scales polynomially in circuit depth while exhaustive classical search scales exponentially.

\subsubsection{Why the Unoptimized Baseline Scores Lower}
Despite 100\% detail retention, the baseline achieves only 50\% voice score and 6.2\% guardrail strength. In observed calls, this produced TTS verbalization of markdown symbols---extending utterance duration by 2.1 seconds on average---and frequent multi-sentence preambles. The QUBO-selected prompt eliminates both failure modes.

\subsubsection{Random Templates Are Harmful}
Random selection (row 2) reduces fitness by 68.5\% and detail retention to 45.1\%, because generic template slots (\texttt{[Agent Name]}, \texttt{[primary purpose]}) are not consistently filled. This validates the need for structured optimization even with template-based prompting.

% =========================================================================
% V. FUTURE SCOPE
% =========================================================================
\section{Future Scope}

\subsection{Modular Prompt Assembly at Scale}
The current system selects among $n = 4$ monolithic variants. With $m$ independent instruction modules, the search space is $2^m$. The QUBO Hamiltonian scales to this regime without modification; at $m \approx 25$, classical greedy becomes intractable while quantum circuit depth grows linearly.

\subsection{Physical QPU Deployment}
Statevector simulation provides exact ground-state results. Physical NISQ hardware introduces gate infidelity. Future work will apply Zero-Noise Extrapolation~\cite{temme2017} and Pauli twirling for fidelity preservation.

\subsection{Dynamic Mid-Call Prompt Morphing}
Tracking caller sentiment over a sliding dialogue window allows real-time Hamiltonian updates:
\begin{equation}
\mathbf{x}^*(t) = \arg\min_{\mathbf{x}} \mathcal{H}_C(\mathbf{x};\, \boldsymbol{\theta}_{\text{caller}}(t))
\end{equation}
enabling real-time persona switching between operational, de-escalation, and empathetic modes.

% =========================================================================
% VI. CONCLUSION
% =========================================================================
\section{Conclusion}

This paper presented a production-grade conversational voice agent with Quantum Prompt Optimization (QPO). The QUBO/QAOA formulation maps five measurable voice-agent quality criteria to an Ising spin Hamiltonian, selecting the optimal system prompt among generated candidates. The selected prompt improved voice appropriateness from 50\% to 90\% and guardrail strength from 6.2\% to 18.8\% over the unoptimized baseline.

We report honestly that for $n = 4$ candidates, Classical Greedy reaches the same optimum in $<$0.01\,ms. The quantum advantage materializes in the modular prompt assembly regime, where quadratic constraint encoding prevents conflicting module co-selection---a capability that classical linear scoring cannot express. Quantum optimization contributes \textbf{0\,ms} to call latency. The full runtime pipeline achieves a median end-to-end latency of \textbf{347.5\,ms}, with STT and LLM TTFT constituting 81\% of total latency.

% =========================================================================
% REFERENCES
% =========================================================================
\begin{thebibliography}{00}
\bibitem{mctear2022} M. McTear, \emph{Conversational AI: Dialogue Systems, Conversational Agents, and Chatbots}. Cham: Springer, 2022.
\bibitem{larson2020} J. Larson and D. R. Traum, ``Turn-taking and latency in spoken dialogue systems,'' \emph{IEEE/ACM Trans. Audio, Speech, Lang. Process.}, vol.~28, pp.~2911--2923, 2020.
\bibitem{baevski2020} A. Baevski \emph{et al.}, ``wav2vec 2.0,'' in \emph{Proc. NeurIPS}, vol.~33, 2020, pp.~12449--12460.
\bibitem{lewis2020} P. Lewis \emph{et al.}, ``Retrieval-augmented generation for knowledge-intensive NLP tasks,'' in \emph{Proc. NeurIPS}, 2020, pp.~9459--9474.
\bibitem{robertson2009} S. Robertson and H. Zaragoza, ``The probabilistic relevance framework: BM25 and beyond,'' \emph{Found. Trends Inf. Retr.}, vol.~3, no.~4, pp.~333--389, 2009.
\bibitem{date2023} P. Date \emph{et al.}, ``Combinatorial optimization on quantum computers,'' \emph{IEEE Trans. Quantum Eng.}, vol.~4, pp.~1--22, 2023.
\bibitem{farhi2014} E. Farhi, J. Goldstone, and S. Gutmann, ``A quantum approximate optimization algorithm,'' \emph{arXiv:1411.4028}, 2014.
\bibitem{glover2019} F. Glover, G. Kochenberger, and Y. Du, ``Quantum bridge analytics I: QUBO models,'' \emph{4OR}, vol.~17, no.~4, pp.~335--371, 2019.
\bibitem{temme2017} K. Temme, S. Bravyi, and J. M. Gambetta, ``Error mitigation for short-depth quantum circuits,'' \emph{Phys. Rev. Lett.}, vol.~119, 2017.
\end{thebibliography}

\end{document}
